\begin{thebibliography}{99}
\small
\linespread{0.94}\selectfont
\setlength{\itemsep}{0pt}
\ifcsname pubref@DeDekenLowen\endcsname
\bibitem{DeDekenLowen} O. De Deken, W. Lowen,
\emph{Abelian and derived deformations in the presence of $\mathbb Z$-generating geometric helices},
J. Noncommut. Geom. 5 (2011), 477--505,
\href{https://doi.org/10.4171/JNCG/83}{doi:10.4171/JNCG/83}, Theorem 4.5, Proposition 4.6, pp. 495--496.
\fi
\ifcsname pubref@StacksPi1\endcsname
\bibitem{StacksPi1} The Stacks Project Authors,
\emph{The Stacks Project},
\href{https://stacks.math.columbia.edu/tag/0BTX}{Tag 0BTX}, Lemma 58.14.3
(arithmetic--geometric fundamental group exact sequence).
\fi
\ifcsname pubref@StacksPerfect\endcsname
\bibitem{StacksPerfect} The Stacks Project Authors,
\emph{The Stacks Project},
\href{https://stacks.math.columbia.edu/tag/0BDK}{Tag 0BDK}, Lemma 36.31.3.
\fi
\ifcsname pubref@EG\endcsname
\bibitem{EG} P. Etingof, V. Ginzburg,
\emph{Noncommutative del Pezzo surfaces and Calabi--Yau algebras},
J. Eur. Math. Soc. 12 (2010), 1371--1416.
\href{https://arxiv.org/abs/0709.3593v4}{arXiv:0709.3593v4}.
Theorem 3.3.2(i),(ii), Proposition 7.3.1, Appendix 9.1。
\fi
\ifcsname pubref@LBS\endcsname
\bibitem{LBS} L. Le Bruyn, S. P. Smith, M. Van den Bergh,
\emph{Central extensions of three dimensional Artin--Schelter regular algebras},
Math. Z. 222 (1996), 171--212, Theorem 2.6, Corollary 2.7。
\fi
\ifcsname pubref@BG\endcsname
\bibitem{BG} A. Braverman, D. Gaitsgory,
\emph{Poincar\'e--Birkhoff--Witt theorem for quadratic algebras of Koszul type},
J. Algebra 181 (1996), 315--328, Theorem 0.5。
\fi
\ifcsname pubref@IS\endcsname
\bibitem{IS} N. Iyudu, S. Shkarin,
\emph{Three dimensional Sklyanin algebras and Gr\"obner bases},
\href{https://arxiv.org/abs/1601.00564}{arXiv:1601.00564}, Theorem 1.1, Corollary 6.1。
\fi
\ifcsname pubref@NS\endcsname
\bibitem{NS} C. Negron, T. Schedler,
\emph{The Hochschild cohomology ring of a global quotient orbifold},
Adv. Math. 364 (2020), 106978.
\href{https://arxiv.org/abs/1809.08715}{arXiv:1809.08715}, Theorem 1.1, Appendix A。
\fi
\ifcsname pubref@FIOS\endcsname
\bibitem{FIOS} E. Faber, C. Ingalls, S. Okawa, M. Satriano,
\emph{On stacky surfaces and noncommutative surfaces},
\href{https://arxiv.org/abs/2206.13359v2}{arXiv:2206.13359v2}, 8 February 2024,
Theorem 1.5, Corollary 1.7。Brauer捩れ付きの比較として使用する。
\fi
\ifcsname pubref@FPSS\endcsname
\bibitem{FPSS} L. Fu, M. Porta, S. Scherotzke, N. Sibilla,
\emph{Hochschild--Kostant--Rosenberg isomorphism for derived Deligne--Mumford stacks},
\href{https://arxiv.org/abs/2509.00501v3}{arXiv:2509.00501v3}, 20 January 2026,
Proposition 6.17。一般root stackの恒等部分への帰着を比較し、$n=2$ に限る後半の式を $n=3$ へ流用しない。
\fi
\ifcsname pubref@TerwilligerDAHA\endcsname
\bibitem{TerwilligerDAHA} P. Terwilliger,
\emph{The universal Askey--Wilson algebra and DAHA of type $(C_1^\vee,C_1)$},
SIGMA 9 (2013), 047, \href{https://doi.org/10.3842/SIGMA.2013.047}{doi:10.3842/SIGMA.2013.047},
Definitions 2.1, 2.2, 2.6; Lemmas 2.3, 2.5, 2.8.
\fi
\ifcsname pubref@GoldmanTrace\endcsname
\bibitem{GoldmanTrace} W. M. Goldman,
\emph{Trace coordinates on Fricke spaces of some simple hyperbolic surfaces},
in \emph{Handbook of Teichm\"uller Theory}, Vol. II (2009), 611--684,
\href{https://doi.org/10.4171/055-1/16}{doi:10.4171/055-1/16},
\S5.2.1, equation (5.2.1.2).
\fi
\ifcsname pubref@GoldmanToledo\endcsname
\bibitem{GoldmanToledo} W. M. Goldman, D. Toledo,
\emph{Affine cubic surfaces and relative $\mathrm{SL}(2)$-character varieties of compact surfaces},
\href{https://arxiv.org/abs/1006.3838}{arXiv:1006.3838}.
Theorem 1: every smooth-at-infinity cubic with generic tritangent plane is a four-holed-sphere character variety.
\fi
\ifcsname pubref@SchockD4\endcsname
\bibitem{SchockD4} N. Schock,
\emph{The $W(E_6)$-invariant birational geometry of the moduli space of marked cubic surfaces},
Math. Nachr. 297 (2024), 2639--2666,
\href{https://doi.org/10.1002/mana.202300459}{doi:10.1002/mana.202300459}.
Section 2: tritangent triangle and its orthogonal $D_4$ root lattice.
\fi
\ifcsname pubref@IISDiscriminant\endcsname
\bibitem{IISDiscriminant} M. Inaba, K. Iwasaki, M.-H. Saito,
\emph{Dynamics of the sixth Painlev\'e equation},
S\'eminaires \& Congr\`es 14 (2006), 103--167,
\href{https://emis.de/ft/47232}{original text}, \S9.1, Definition 9.1.
Fricke判別式をトレースパラメータへ引き戻した場合の二乗因子について参照する。
本稿の$N_5^2M_5$との同一視は主張しない。
\fi
\ifcsname pubref@HavlicekPosta\endcsname
\bibitem{HavlicekPosta} M. Havl\'{\i}\v{c}ek, S. Po\v{s}ta,
\emph{Center of quantum algebra $U'_q(\mathrm{so}_3)$},
J. Math. Phys. 52 (2011), 043521,
\href{https://doi.org/10.1063/1.3579992}{doi:10.1063/1.3579992}.
三乗根を含むroot-of-unity中心の構造について参照し、規約の完全な同一視は追加検証事項。
\fi
\ifcsname pubref@CantatLorayCayley\endcsname
\bibitem{CantatLorayCayley} S. Cantat, F. Loray,
\emph{Dynamics on Character Varieties and Malgrange irreducibility of Painlev\'e VI equation},
Ann. Inst. Fourier 59 (2009), 2927--2978,
\href{https://doi.org/10.5802/aif.2512}{doi:10.5802/aif.2512}, \S1.5.
\fi
\ifcsname pubref@JangBoundary\endcsname
\bibitem{JangBoundary} S. u. Jang,
\emph{Tropical Dynamics of Markov Surfaces},
\href{https://arxiv.org/abs/2306.11357v2}{arXiv:2306.11357v2},
Lemma 4.8 and Theorem 4.10. 取得した2025年12月版の題名と番号を使用する。
\fi
\ifcsname pubref@MikoszWeberTriality\endcsname
\bibitem{MikoszWeberTriality} M. Mikosz, A. Weber,
\emph{Triality, characteristic classes, $D_4$ and $G_2$ singularities},
J. Homotopy Relat. Struct. 10 (2015), 995--1011,
\href{https://doi.org/10.1007/s40062-014-0094-3}{doi:10.1007/s40062-014-0094-3}, Proposition 3, \S5.
\fi
\ifcsname pubref@DixonElliptope\endcsname
\bibitem{DixonElliptope} K. Dixon,
\emph{The multi-moment maps of the nearly K\"ahler $S^3\times S^3$},
Geom. Dedicata 200 (2019), 351--362,
\href{https://doi.org/10.1007/s10711-018-0375-9}{doi:10.1007/s10711-018-0375-9}, Lemma 3.5.
\fi
\ifcsname pubref@HikamiCluster\endcsname
\bibitem{HikamiCluster} K. Hikami,
\emph{Note on character varieties and cluster algebras},
SIGMA 15 (2019), 003,
\href{https://doi.org/10.3842/SIGMA.2019.003}{doi:10.3842/SIGMA.2019.003}.
\fi
\ifcsname pubref@AWAvatars\endcsname
\bibitem{AWAvatars} N. Cramp\'e et al.,
\emph{The Askey--Wilson algebra and its avatars},
J. Phys. A 54 (2021), 063001,
\href{https://doi.org/10.1088/1751-8121/abd783}{doi:10.1088/1751-8121/abd783}.
\fi
\ifcsname pubref@CookeLacabanne\endcsname
\bibitem{CookeLacabanne} J. Cooke, N. Lacabanne,
\emph{Higher rank Askey--Wilson algebras as skein algebras},
Ann. Inst. Fourier 76 (2026), 1055--1117,
\href{https://doi.org/10.5802/aif.3729}{doi:10.5802/aif.3729}.
\fi
\ifcsname pubref@BousseauSkein\endcsname
\bibitem{BousseauSkein} P. Bousseau,
\emph{Strong positivity for the skein algebras of the $4$-punctured sphere and the $1$-punctured torus},
Comm. Math. Phys. 398 (2023), 1--58,
\href{https://doi.org/10.1007/s00220-022-04512-9}{doi:10.1007/s00220-022-04512-9}.
\fi
\ifcsname pubref@EOR\endcsname
\bibitem{EOR} P. Etingof, A. Oblomkov, E. Rains,
\emph{Generalized double affine Hecke algebras of rank 1 and quantized del Pezzo surfaces},
Adv. Math. 212 (2007), 749--796,
\href{https://doi.org/10.1016/j.aim.2006.11.008}{doi:10.1016/j.aim.2006.11.008}.
\fi
\ifcsname pubref@CMR\endcsname
\bibitem{CMR} L. Chekhov, M. Mazzocco, V. Rubtsov,
\emph{Quantised Painlev\'e monodromy manifolds, Sklyanin and Calabi--Yau algebras},
Adv. Math. 376 (2021), 107442,
\href{https://doi.org/10.1016/j.aim.2020.107442}{doi:10.1016/j.aim.2020.107442},
\href{https://arxiv.org/abs/1905.02772}{arXiv:1905.02772}.
原著v2のDefinition 1.1、式(1.5)、Theorem 1.5、Proposition 4.2を参照。
\fi
\ifcsname pubref@Rains\endcsname
\bibitem{Rains} E. M. Rains,
\emph{Filtered deformations of elliptic algebras},
\href{https://arxiv.org/abs/2107.13540v2}{arXiv:2107.13540v2}, 17 February 2022,
Propositions 2.2, 2.3, Example 2.1, Section 6末尾のConjecture 1とRemark。
\fi
\ifcsname pubref@LVdB\endcsname
\bibitem{LVdB} W. Lowen, M. Van den Bergh,
\emph{Hochschild cohomology of abelian categories and ringed spaces},
Adv. Math. 198 (2005), 172--221,
\href{https://arxiv.org/abs/math/0405227}{arXiv:math/0405227}。
\fi
\ifcsname pubref@LowenObject\endcsname
\bibitem{LowenObject} W. Lowen,
\emph{Hochschild cohomology, the characteristic morphism and derived deformations},
Compos. Math. 144 (2008), 1557--1580,
\href{https://arxiv.org/abs/0707.2602}{arXiv:0707.2602}。
\fi
\ifcsname pubref@Repo\endcsname
\bibitem{Repo} \emph{nc-cubic-moduli research repository},
\href{https://github.com/uedakazushi/nc-cubic-moduli}{GitHub main}。
本稿の依存する個別ファイル・ラベル・確認済み状態は研究ノートの入力表に固定する。
\fi
\ifcsname pubref@CompanionE6\endcsname
\bibitem{CompanionE6}\emph{三乗根量子三次曲面の正常形と形式変形},
日本語研究原稿、2026年11月18日、\PublicationVersion。
\href{https://github.com/uedakazushi/nc-cubic-e6/blob/main/papers/nc-cubic-e6.tex}{\texttt{nc-cubic-e6.tex}}。
\fi
\ifcsname pubref@CompanionFricke\endcsname
\bibitem{CompanionFricke}\emph{量子三次曲面のFricke幾何と周期},
日本語研究原稿、2026年11月18日、\PublicationVersion。
\href{https://github.com/uedakazushi/nc-cubic-e6/blob/main/papers/fricke-geometry.tex}{\texttt{fricke-geometry.tex}}。
\fi
\ifcsname pubref@CompanionBoundary\endcsname
\bibitem{CompanionBoundary}\emph{三次曲面のPell力学と境界特異点},
日本語研究原稿、2026年11月18日、\PublicationVersion。
\href{https://github.com/uedakazushi/nc-cubic-e6/blob/main/papers/cubic-boundary-prym.tex}{\texttt{cubic-boundary-prym.tex}}。
\fi
\ifcsname pubref@AOUCubic\endcsname
\bibitem{AOUCubic}
T. Abdelgadir, S. Okawa, K. Ueda,
\emph{Compact moduli of marked noncommutative cubic surfaces},
\href{https://arxiv.org/abs/2404.00175v1}{arXiv:2404.00175v1} (2024).
\S9, equations (9.1)--(9.10), Lemma 9.2, Theorem 9.4.
\fi
\ifcsname pubref@OUASSurfaces\endcsname
\bibitem{OUASSurfaces}
S. Okawa, K. Ueda,
\emph{Marked Artin--Schelter surfaces of del Pezzo types},
\href{https://arxiv.org/abs/2508.04986v1}{arXiv:2508.04986v1} (2025).
Definition 1.1.
\fi
\ifcsname pubref@BoundaryDHKK\endcsname
\bibitem{BoundaryDHKK} G. Dimitrov, F. Haiden, L. Katzarkov, M. Kontsevich,
\emph{Dynamical systems and categories},
\href{https://arxiv.org/abs/1307.8418}{arXiv:1307.8418},
Definitions 2.1, 2.4, Proposition 2.2, Lemma 2.3, Theorem 2.6.
\fi
\ifcsname pubref@BoundaryGoldmanNeumann\endcsname
\bibitem{BoundaryGoldmanNeumann} W. M. Goldman, W. D. Neumann,
\emph{Homological action of the modular group on some cubic moduli spaces},
Math. Res. Lett. 12 (2005), 575--591,
\href{https://doi.org/10.4310/MRL.2005.v12.n4.a11}{doi:10.4310/MRL.2005.v12.n4.a11},
\href{https://arxiv.org/abs/math/0402039}{arXiv:math/0402039},
Theorems 1--2, Proposition 3.
\fi
\ifcsname pubref@BoundaryDyckerhoff\endcsname
\bibitem{BoundaryDyckerhoff} T. Dyckerhoff,
\emph{Compact generators in categories of matrix factorizations},
Duke Math. J. 159 (2011), 223--274,
\href{https://doi.org/10.1215/00127094-1415869}{doi:10.1215/00127094-1415869},
\href{https://arxiv.org/abs/0904.4713v5}{arXiv:0904.4713v5},
Theorem 5.9, Corollary 6.5.
\fi
\ifcsname pubref@BoundaryHuaKeller\endcsname
\bibitem{BoundaryHuaKeller} Z. Hua, B. Keller,
\emph{Cluster categories and rational curves},
Geom. Topol. 28 (2024), 2569--2634,
\href{https://doi.org/10.2140/gt.2024.28.2569}{doi:10.2140/gt.2024.28.2569},
\href{https://arxiv.org/abs/1810.00749v6}{arXiv:1810.00749v6},
Theorem 5.9, pp.36--37.
\fi
\ifcsname pubref@BoundaryMatherYau\endcsname
\bibitem{BoundaryMatherYau} J. N. Mather, S. S.-T. Yau,
\emph{Classification of isolated hypersurface singularities by their moduli algebras},
Invent. Math. 69 (1982), 243--251,
\href{https://doi.org/10.1007/BF01399504}{doi:10.1007/BF01399504}.
\fi
\ifcsname pubref@BoundaryMumford\endcsname
\bibitem{BoundaryMumford} D. Mumford,
\emph{Prym Varieties I}, in \emph{Contributions to Analysis},
Academic Press (1974), 325--350,
\href{https://doi.org/10.1016/B978-0-12-044850-0.50032-0}{doi:10.1016/B978-0-12-044850-0.50032-0},
\S\S1--3.
\fi
\ifcsname pubref@BoundaryMirandaPersson\endcsname
\bibitem{BoundaryMirandaPersson} R. Miranda, U. Persson,
\emph{On extremal rational elliptic surfaces},
Math. Z. 193 (1986), 537--558,
\href{https://doi.org/10.1007/BF01160474}{doi:10.1007/BF01160474},
Table 5.2.
\fi
\ifcsname pubref@BoundaryHowald\endcsname
\bibitem{BoundaryHowald} J. Howald,
\emph{Multiplier Ideals of Sufficiently General Polynomials},
\href{https://arxiv.org/abs/math/0303203}{arXiv:math/0303203},
Theorems 11--12, Corollary 13.
\fi
\ifcsname pubref@BoundaryCasals\endcsname
\bibitem{BoundaryCasals} R. Casals,
\emph{Lagrangian skeleta and plane curve singularities},
J. Fixed Point Theory Appl. 24 (2022), 34,
\href{https://doi.org/10.1007/s11784-022-00939-8}{doi:10.1007/s11784-022-00939-8},
Theorem 1.1, Corollary 1.2, \S2.1.
\fi
\ifcsname pubref@BoundaryCamara\endcsname
\bibitem{BoundaryCamara} L. M. C\^amara,
\emph{On the classification of quasi-homogeneous curves},
\href{https://arxiv.org/abs/1009.1664v1}{arXiv:1009.1664v1},
Theorem 1(iii), \S4.3.
\fi
\ifcsname pubref@FrickeFriedman\endcsname
\bibitem{FrickeFriedman}
R. Friedman, \emph{On the geometry of anticanonical pairs},
\href{https://arxiv.org/abs/1502.02560v2}{arXiv:1502.02560v2} (2016),
Proposition 3.12.
\fi
\ifcsname pubref@FrickeBoalchPaluba\endcsname
\bibitem{FrickeBoalchPaluba}
P. Boalch, R. Paluba, \emph{Symmetric cubic surfaces and $G_2$ character varieties},
J. Algebraic Geom. 25 (2016), 607--631,
\href{https://doi.org/10.1090/jag/668}{doi:10.1090/jag/668},
Lemma 6, Corollary 11, Proposition 18, Remark 19.
\fi
\ifcsname pubref@PymR13\endcsname
\bibitem{PymR13} B. Pym,
\emph{Quantum deformations of projective three-space},
Adv. Math. 281 (2015), 1216--1241,
\href{https://doi.org/10.1016/j.aim.2015.06.005}{doi:10.1016/j.aim.2015.06.005},
\href{https://arxiv.org/abs/1403.6444}{arXiv:1403.6444}, \S3.4, (9).
\fi
\ifcsname pubref@CompanionAmbient\endcsname
\bibitem{CompanionAmbient}
\emph{非可換三次曲面のambient射と$R(1,3)$成分},
\href{https://github.com/uedakazushi/nc-cubic-ambient/blob/main/papers/nc-cubic-ambient.tex}
{nc-cubic-ambient research manuscript},
\texttt{thm:r13-affine}, \texttt{prop:center}, \texttt{sec:surface-fiber}.
引用範囲は次数付き中心拡大の同型商と中心三次ペンシルであり、
AOUからの支配性の追加仮定は用いない。
\fi
\end{thebibliography}
